\styledtitle{Conclusion générale}
\begin{spacing}{1.5}
\vspace{1cm}
\label{chap:ConclusionPers}
\addstarredchapter{Conclusion générale}
\lhead{Conclusion générale et perspectives}
\rhead{\thepage}

Ce projet de fin d'études a permis de concevoir et de développer un système
intelligent capable de détecter les falsifications des mesures de température
transmises par les capteurs d'une chaîne du froid pharmaceutique. En combinant
la construction d'un jeu de données avec attaques injectées, la modélisation
prédictive, l'explicabilité des alertes et la visualisation interactive, nous
avons réalisé une application web, baptisée VigiloChain, qui constitue un outil
d'aide à la décision pour les opérateurs chargés de protéger les médicaments
sensibles.

Les résultats obtenus démontrent la pertinence de la démarche adoptée. Sur le
jeu de test, jamais utilisé pour choisir le modèle ni le seuil, le modèle retenu
atteint une précision de 92,7\,\% et un rappel de 65,2\,\%, soit un F1 de 0,765.
Au niveau des épisodes, 70 attaques sur 73 sont détectées, avec un délai médian
de 4 minutes et 0,35 fausse alerte par heure. Les performances proches entre
validation et test indiquent l'absence de sur-apprentissage marqué. L'intégration
d'outils modernes tels que Extra Trees, SHAP, FastAPI et React a permis d'obtenir
un détecteur performant et de rendre ses résultats accessibles de manière claire :
chaque alerte est expliquée par les variables qui l'ont provoquée, puis
accompagnée d'un rapport de recommandations pour l'opérateur.

Il est à noter, en outre, que les attaques utilisées dans cette étude sont
injectées dans des traces réelles avec des signatures connues, et que le temps
réel du tableau de bord repose sur le rejeu de mesures enregistrées. Les
performances présentées ne reflètent donc pas encore celles d'un déploiement sur
des capteurs en service. Par conséquent, il s'avère essentiel, pour tout travail
futur, de valider le système sur des flux réels, avec de véritables tentatives de
falsification ou des essais en laboratoire, afin de confirmer la pertinence des
résultats obtenus.

Dans le cadre de la recherche d'axes d'amélioration, l'extension du jeu de
données est une première piste : elle consisterait à ajouter de nouveaux
scénarios d'attaque, d'autres capteurs et d'autres types de produits.
Notre application intègre déjà un mécanisme de rejeu qui fait passer les
mesures dans le modèle sauvegardé avec le seuil et la règle de persistance
retenus, et qui lève des alertes comme le ferait un flux réel. Une perspective
d'évolution consiste à connecter directement l'application aux capteurs, par
exemple au moyen d'un protocole de messagerie de l'Internet des objets comme MQTT.
Cette automatisation permettrait de collecter et d'analyser en continu les
mesures réelles, sans passer par un fichier. De plus, l'extension du système à
plusieurs sites et à plusieurs chaînes du froid, avec une supervision centralisée
et des droits d'accès par site, permettrait de dépasser le cadre d'un seul
entrepôt.

Une autre perspective intéressante consisterait à remplacer l'hypothèse par un
diagnostic. Le modèle actuel est binaire : il répond seulement « attaque » ou
« normal », et la famille d'attaque proposée dans le rapport est déduite des
variables qui ont le plus pesé dans l'alerte. Un modèle multiclasse, entraîné à
prédire directement la famille de manipulation, offrirait à l'opérateur une
information plus directe et plus fiable. Bien que cette tâche n'ait pas pu être
réalisée dans le cadre de ce projet par manque de temps, elle constitue une piste
prometteuse pour enrichir la portée de notre solution.

Enfin, l'adaptation de la gravité au produit transporté représenterait une
amélioration notable. Les seuils de gravité actuels sont des choix de conception,
indépendants du type de médicament et de l'écart réel à la plage de conservation.
En les reliant aux exigences de chaque produit, et en ajoutant une notification
de l'opérateur et un suivi de chaque alerte jusqu'à sa clôture, les rapports
générés deviendraient plus précis et directement exploitables dans le processus
de décision.

Sur le plan personnel, ce projet a représenté une expérience extrêmement
formatrice. Nous avons approfondi nos compétences en science des données
(construction de variables causales, évaluation rigoureuse, explicabilité),
mais aussi développé de nouvelles aptitudes en développement web et en
visualisation interactive. Nous avons appris à mener un projet de bout en bout,
depuis la définition du besoin métier jusqu'au déploiement d'une solution
concrète et testée, en suivant la démarche CRISP-DM. Ce travail nous a permis de
renforcer notre esprit analytique, notre autonomie et notre rigueur scientifique.

En conclusion, ce projet constitue une étape importante dans notre évolution
professionnelle, en nous donnant la confiance et l'expérience nécessaires pour
aborder de futurs défis, tout en contribuant à apporter des solutions concrètes
à un problème réel : la protection de la chaîne du froid pharmaceutique.

\end{spacing}
